\subsubsection{Características do Experimento Fatorial}

O experimento foi estruturado como um fatorial completo $2 \times 3$ com os parâmetros a seguir:

\begin{align*}
\text{Níveis de Temperatura} &= 2 \\
\text{Níveis de Pressão} &= 3 \\
\text{Total de Combinações} &= 6 \\
\text{Total de Observações} &= 30
\end{align*}

\paragraph{Modelo Fatorial}
O modelo estatístico é dado por:

\[Y_{ijk} = \mu + \tau_i + \gamma_j + (\tau\gamma)_{ij} + \epsilon_{ijk}\]

onde: $Y_{ijk}$ é a resposta; $\tau_i$ é o efeito da Temperatura $i$; $\gamma_j$ é o efeito da Pressão $j$; $(\tau\gamma)_{ij}$ é o efeito de interação; $\epsilon_{ijk}$ é o erro aleatório.

\paragraph{Conceito de Interação}
A interação ocorre quando o efeito de um fator depende do nível do outro fator. Quando presente (p-valor $< 0.05$), a interação significa que:

\begin{itemize}
\item Não é possível fazer recomendações sobre um fator independentemente do outro
\item A análise deve focar nas combinações ótimas de níveis dos fatores
\item Os gráficos de perfil mostram linhas não-paralelas
\end{itemize}

\paragraph{Interpretação dos Resultados Esperados}
Para este experimento, espera-se:

\begin{enumerate}
\item \textbf{Efeito Principal de Temperatura:} Significativo (efeito esperado)
\item \textbf{Efeito Principal de Pressão:} Significativo (efeito esperado)
\item \textbf{Interação Temperatura $\times$ Pressão:} Muito significativa ($p \ll 0.05$)
\item \textbf{Conclusão:} A pressão ótima depende criticamente da temperatura escolhida
\end{enumerate}

\paragraph{Recomendação Prática}
A existência de interação exige que a recomendação se baseie na combinação ótima de Temperatura $\times$ Pressão, não em fatores isolados.

